\documentclass[11pt]{article}

% --- packages (all standard; compiles with a vanilla TeX Live install) ---
\usepackage[margin=1in]{geometry}
\usepackage{amsmath}
\usepackage{booktabs}
\usepackage{tabularx}
\usepackage{array}
\usepackage{enumitem}
\usepackage{xcolor}
\usepackage[hidelinks]{hyperref}

% Print field/code names in a monospace style. Underscores are escaped inline.
\newcommand{\code}[1]{\texttt{#1}}

% A self-contained figure placeholder: a framed caption-box so the document
% compiles with no external image files. Replace each with \includegraphics
% once the real plots exist.
\newcommand{\figplaceholder}[1]{%
  \fbox{\parbox[c][3.2cm][c]{0.92\linewidth}{\centering\itshape #1}}%
}

\newcolumntype{L}{>{\raggedright\arraybackslash}X}

\title{Weld Agent Study --- Data Analysis Plan\\[2pt]
\large How the analysis maps onto the paper}
\author{Andrew Sadler (Student Investigator) \and Dr.\ Jie Wu (PI)\\
Computer Science, Michigan Technological University}

\begin{document}
\maketitle

\begin{abstract}
\noindent
This document explains how the study's data becomes the Results section of the
paper: the metrics we compute, the statistical models we fit, the figures we
report, and the order in which the argument is presented. Every measure is tied
to the field that produces it (a telemetry event type and payload key, or a
survey \code{item\_id}) so the analysis is reproducible from the event log and
the survey export. It is a plan, not yet results: numbers and figures are
placeholders to be filled after data collection.
\end{abstract}

\section{The spine of the argument}\label{sec:spine}

The whole analysis ladders up to one claim: \emph{unifying task, communication,
and code context and making it accessible to an AI agent helps---and the benefit
appears exactly where it should (cross-surface work), not everywhere.} Showing
where the effect is \emph{absent} (comprehension and single-surface control
tasks) is what makes the place it \emph{is} present credible.

The independent variable is a three-rung \emph{agent-scope ladder}:
\begin{itemize}[nosep]
  \item \textbf{C1} --- agent sees chat only, and advises only;
  \item \textbf{C2} --- agent sees chat $+$ board, and can act on both;
  \item \textbf{C3} --- agent sees chat $+$ board $+$ GitHub, and can act across all three.
\end{itemize}
Moving up the ladder \emph{is} increasing integration, so ``does integration
help?'' becomes the concrete, testable ``does $\text{C3} > \text{C2} >
\text{C1}$, and on which tasks?''

\paragraph{Scope, not context breadth alone --- and why the naming matters.}
Visibility and capability rise together across the ladder: an agent cannot act on
a board it cannot see, so the two are coupled by construction rather than by
oversight. We therefore name the manipulated variable \emph{agent scope}
throughout, and not ``context breadth,'' which would imply capability was held
fixed. The consequence for interpretation is specific and we state it up front:
\textbf{C3-vs-C1 confounds seeing with doing.} C1 can only talk; C3 can act
across three surfaces. On cross-surface tasks that contrast can hardly fail, and
its most parsimonious explanation is that an agent which performs the steps saves
the user from performing them --- a real and worthwhile finding about agent
scope, but not evidence that \emph{context} is what did the work. The contrast
that isolates context is \textbf{C2-vs-C3}: both act, and they differ only in
whether GitHub is in scope. At $\sim$7--10 per cell that comparison is
underpowered, and we report it as an estimate with its interval rather than a
test. Read together: C3-vs-C1 establishes that agent scope matters and functions
largely as a manipulation check; C2-vs-C3 is the theoretically clean but
imprecise estimate of the GitHub rung; and the study's distinctive contribution
is H5, where the confound does not apply at all, since trust is compared against
each condition's \emph{own} measured effectiveness.

\subsection*{Research questions and hypotheses}

The study manipulates one variable---how much of the workspace the in-app agent
can see and act on (the C1$\rightarrow$C2$\rightarrow$C3 scope ladder)---measured
against each participant's own no-agent baseline (C0). The questions stay inside
what that supports: the agent's scope over the workspace and how users work with
it, not whether an all-in-one workspace beats separate tools (there is no
fragmented-tools arm to test that), and not context breadth with capability held
constant (the two rise together by construction).

\begin{description}[leftmargin=2.6em,style=nextline,nosep]
  \item[RQ1] \emph{(primary)} Does widening the agent's scope over the workspace
  (chat $\rightarrow$ $+$board $\rightarrow$ $+$GitHub) make the user more
  effective, and on which tasks?
  \item[RQ2] How much do users delegate to vs.\ verify the agent, and does that
  shift with its scope?
  \item[RQ3] Does trust in the agent track its measured effectiveness, or diverge?
\end{description}

\begin{description}[leftmargin=3.0em,style=nextline,nosep]
  \item[H1] (RQ1) Wider scope yields shorter baseline-adjusted agent-task times
  \emph{across conditions}. (Stated between conditions, never as an absolute gain
  over C0: phase order is fixed, so an agent$-$baseline delta also contains
  ordinary practice --- see \S\ref{sec:timing}.)
  \item[H2] (RQ1) The gain is \emph{task-dependent}: minimal on comprehension and
  single-surface tasks, substantial on cross-surface tasks.
  \item[H3] (RQ1) Wider scope lowers self-reported cognitive load relative to
  baseline more than narrow scope.
  \item[H4] (RQ2) Delegation rises with scope, but reliance is
  imperfectly calibrated to where the agent is actually correct.
  \item[H5] (RQ3) Trust and perceived speedup need not track the agent's measured
  effectiveness (the perception--behavior gap).
\end{description}

\section{The design trick: baseline-normalized time}

Each participant first completes six manual tasks (B1--B6, condition C0, no
agent), establishing their own solo speed on each task type
(B1 single-surface; B2--B5 cross-surface; B6 comprehension). Agent-task times are
then expressed as a \emph{ratio to that participant's own type-matched baseline}.
This is a hybrid within-/between-subjects design: condition is between subjects,
but the C0 anchor is within subjects.

Why it matters: people differ enormously in raw speed, and at $N\approx20$--$30$
that person-to-person variance would otherwise swamp the condition effect.
Normalizing to each person's baseline cancels ``this person is simply fast'' and
isolates ``how much did the agent help \emph{this} person.'' It is the single
biggest source of statistical power in a modest sample. The six-task baseline
(vs.\ a single warm-up) buys two things: a type-matched anchor for every agent
task type, and a stabler per-person estimate (averaging the four cross-surface
baselines) for the noisiest cell.

\section{Measures (operationalized)}

\begin{table}[h]
\centering\small
\begin{tabularx}{\linewidth}{@{}L L@{}}
\toprule
\textbf{Construct} & \textbf{Source field(s)} \\
\midrule
Task completion time (\emph{primary outcome}) & \code{task\_completed.task\_time\_ms} --- active engaged time (\S\ref{sec:timing}); clock stops before the per-task rating \\
Wall / away time, interruptions & \code{task\_completed.\{wall\_ms, away\_ms, pause\_count, was\_interrupted\}} \\
Time-to-first-action & first surface event after \code{task\_shown} $-$ \code{task\_shown} timestamp (derived) \\
Completion vs.\ abandonment & \code{task\_completed} / \code{task\_skipped}; \code{session\_abandoned} \\
Correctness (0--3) & manual coding of \code{task\_final\_state} (board $+$ message content) $+$ \code{study\_agent\_turns} $+$ DB \\
Per-task effort (within-person load) & \code{task\_rating.effort} (1--7, every task incl.\ baselines) \\
Perceived agent help & \code{task\_rating.agent\_help} (1--7, agent tasks) \\
Global cognitive load (per phase) & 6 NASA-TLX subscales collected \emph{twice}: \code{tlx\_*\_baseline} (post-baseline interstitial, manual phase) and \code{tlx\_*} (exit survey, agent phase) --- the within-person delta is the H3 signal \\
Usability & \code{sus\_01}\,$\ldots$\,\code{sus\_10} \\
Integration perception & \code{int\_findacross, int\_noticed, int\_vs\_separate, int\_reuse, int\_switch\_perceived} \\
Trust / reliance & \code{trust\_reliable, trust\_predictable, trust\_competent, trust\_rely, trust\_context, trust\_again, trust\_verify, trust\_speed\_perceived} \\
Agent effectiveness / cost & \code{agent\_response.\{latency\_ms, response\_type, num\_tool\_calls, prompt\_tokens, completion\_tokens, total\_tokens\}}; \code{agent\_tool\_call.\{allowed, executed, error\}} \\
Delegation / division of labor & \code{agent\_tool\_call} (executed) vs.\ user action events; \code{agent\_suggestion\_accepted/\_rejected}, \code{agent\_output\_edited}, \code{agent\_action\_undone} \\
Cross-reference (integration usage) & \code{card\_referenced\_from\_chat}, \code{pr\_referenced\_from\_card}, \code{pr\_referenced\_from\_chat}, \code{issue\_referenced\_from\_card} \\
View-switch rate & count of \code{view\_changed} $\div$ active time \\
Scoping proof & \code{agent\_context\_sent} snapshot per invocation (\S\ref{sec:manip}) \\
\bottomrule
\end{tabularx}
\caption{Each construct and the data that produces it. \code{chat\_referenced\_from\_card} and \code{card\_referenced\_from\_pr} have no UI affordance and are intentionally never emitted---treat their absence as N/A, not missing data.}
\label{tab:measures}
\end{table}

\paragraph{Derived headline quantity (descriptive).} \emph{Baseline-normalized
time} $=$ agent-task active time $\div$ that participant's own type-matched C0
baseline: comprehension agent tasks (T1, T2) $\div$ B6; single-surface (T3, T4)
$\div$ B1; cross-surface (T5, T6) $\div$ $\operatorname{mean}(\text{B2},
\text{B3},\text{B4},\text{B5})$. Every agent task type thus has a like-for-like
anchor; the cross-surface anchor averages four baselines and is the most stable.
This ratio is the quantity we \emph{plot} (Figs.\ \ref{fig:interaction},
\ref{fig:percep}) because it is intuitive; the \emph{inferential} model below does
not pre-divide---it carries the baseline as a covariate so its measurement error
is propagated rather than assumed away (\S\ref{sec:primary}), which matters here
because the single-surface and comprehension anchors rest on one baseline task
each.

\section{Statistical models}

\subsection{Primary model (RQ1 / H1 / H2)}\label{sec:primary}

A linear mixed-effects model (ANCOVA on the log scale) for the active time of
agent task $j$ from participant $i$, adjusting for that participant's own
type-matched C0 baseline as a covariate:
\begin{equation}
\log(\textit{active\_time}_{ij}) = \beta_0
  + \beta_{\text{base}}\,\log(\textit{baseline}_{ij})
  + \beta_{\text{cond}}\,\textit{condition}_i
  + \beta_{\text{type}}\,\textit{task\_type}_j
  + \beta_{\text{int}}\,(\textit{condition}_i \times \textit{task\_type}_j)
  + u_i + \varepsilon_{ij},
\label{eq:primary}
\end{equation}
where $\textit{baseline}_{ij}$ is the type-matched anchor for that observation
(B6 for comprehension, B1 for single-surface, $\operatorname{mean}(\text{B2--B5})$
for cross-surface) and $u_i \sim \mathcal{N}(0,\sigma_u^2)$ is a per-participant
random intercept. In words: a regression predicting agent-task time that holds
each person's own manual speed constant via the $\log(\textit{baseline})$ term,
with the random intercept $u_i$ letting each person have their own residual
level---this is how the model respects that several tasks come from the same
person (their observations are not independent). That is the ``mixed'' part.

\textbf{Why ANCOVA rather than a pre-divided ratio.} Dividing by the baseline
treats it as measured without error---fine for the cross-surface anchor (an
average of four tasks), but the comprehension and single-surface anchors are each
a \emph{single} baseline task, so a hard ratio would inject that one task's noise
straight into the outcome on exactly the control task types whose nulls carry the
``context breadth doesn't \emph{universally} help'' argument. Carrying
$\log(\textit{baseline})$ as a covariate instead lets the model account for the
anchor's uncertainty and estimate $\beta_{\text{base}}$ from the data rather than
forcing it to $1$. The normalized ratio remains the descriptive/figure quantity;
this is the inferential test.

\paragraph{The confirmatory number comes from the model, not the ratio.}
\code{cross\_surface} is the reference level of \textit{task\_type} and C1 of
\textit{condition}, so the C3-vs-C1 cross-surface contrast is a linear
combination of Eq.~\eqref{eq:primary}'s coefficients and is reported with the
standard error the model gives it. This matters twice. First, it is the same
argument as the paragraph above: having refused to pre-divide by the anchor
because that treats it as measured without error, we cannot then take the
confirmatory test from a $t$-test on those very ratios. Second, the C2-vs-C3
contrast differences two coefficients, so its standard error must come from the
model's covariance matrix ($\operatorname{Var}(a-b) = \operatorname{Var}(a) +
\operatorname{Var}(b) - 2\operatorname{Cov}(a,b)$); differencing two
independently quoted intervals overstates the width on precisely the contrast
that isolates the GitHub rung. The per-participant ratio contrast is still
computed and reported beside it as the descriptive companion --- it is what
Fig.~\ref{fig:interaction} plots --- and material disagreement between the two
is itself reportable.

\textbf{The \code{condition $\times$ task\_type} interaction
($\beta_{\text{int}}$) is the headline result.} An interaction means the effect
of condition \emph{depends on} the task type. H2 predicts condition barely
matters on comprehension/single-surface tasks but matters a lot on cross-surface
tasks---which is precisely an interaction, and a stronger finding than a flat
``C3 is faster.'' Because the outcome is modeled on the log scale, we lead with
the \emph{back-transformed, baseline-adjusted multiplicative effect}---e.g.\ ``on
cross-surface tasks C3 takes $0.7\times$ the time of C1, 95\,\% CI
$[\cdot,\cdot]$''---which is directly interpretable and closely tracks what
Fig.~\ref{fig:interaction} plots; standardized effect sizes are reported
secondarily. We lead with effect sizes and CIs rather than $p$-values given the
sample size.

\paragraph{Confirmatory vs.\ supporting (fixed in advance).} At
$N\approx20$--$30$ the full $3\times3$ interaction is the most data-hungry term
in the model and will be estimated imprecisely. We therefore fix, before any data
are collected, a
single \emph{focused contrast} as the confirmatory primary test---C3 vs.\ C1 on
cross-surface tasks (the rung and task type where H1/H2 predict the largest gap,
sharpened by T6, \S\ref{sec:t6})---and treat the omnibus
\code{condition$\times$task\_type} interaction as supporting evidence reported
alongside it. Everything past H3 (the trust/perception analyses, H5) is labeled
\emph{exploratory / hypothesis-generating}, not confirmatory.

\emph{What that contrast can and cannot establish.} C3-vs-C1 varies scope on both
axes at once --- C1 advises, C3 acts across three surfaces --- so a gap there is
evidence that widening agent scope helps, and is \emph{not} evidence that the
extra context specifically is what helped (``Scope, not context breadth alone,''
\S\ref{sec:spine}). We choose it
as the confirmatory test because it is the one contrast this sample can estimate
with any precision, and we report it in those terms: as the study's manipulation
working, with the mechanism left to the underpowered C2-vs-C3 estimate and to
T6 (\S\ref{sec:t6}), where the correct PR is identifiable only with GitHub
context. A reader should come away able to distinguish ``scope helped'' (which we
can show) from ``context breadth helped'' (which we can only estimate).

\paragraph{Sensitivity: what this design can detect.} For the confirmatory
contrast (C3 vs.\ C1, two cross-surface agent tasks, eight participants per
cell), the minimum detectable effect at $\alpha=0.05$ two-sided and 80\,\% power
is $d\approx1.1$--$1.4$ across plausible values of the total log-time SD
($\sigma\in[0.4,0.6]$), the within-participant ICC ($\rho\in[0.4,0.6]$), and the
baseline--outcome correlation ($r\in[0,0.7]$). Under central assumptions
($\sigma=0.5$, $\rho=0.5$, $r=0.5$) it is $d=1.20$---a 45\,\% reduction in task
time---with a 20\,000-replication Monte Carlo check recovering 80--83\,\% power
there. Detecting a 20\,\% reduction would require roughly 50 participants per
cell. Two consequences are pre-committed: a non-significant C2-vs-C3 result is
reported as an interval, never as evidence of absence; and because the
between-participant component dominates, more tasks per person is not a remedy
(doubling cross-surface tasks from two to four moves the detectable effect only
from 45\,\% to 42\,\%). Computed by \code{sensitivity.py}; parameters are assumed
pending pilot estimates and are re-run against realised data.

\subsection{Correctness}

An ordinal (cumulative-link) mixed model
\[
\textit{correctness}_{ij} \sim \textit{condition}_i \times \textit{task\_type}_j + (1\mid \textit{participant}_i),
\]
ordinal because the 0--3 score is ranked categories, not a continuous quantity.
\textbf{Pre-specified fallback.} The 0--3 rubric is expected to pile up at 3 on
the control tasks, and with $\sim$7--10 per cell the cumulative-link mixed model
may not converge (sparse cells, ceiling). If it fails to converge or the
proportional-odds assumption is violated, we collapse the score to a binary
\emph{fully correct} ($=3$) vs.\ \emph{not} ($<3$) and report per-cell
proportions with exact (Wilson) confidence intervals, plus a mixed-effects
logistic model where it converges and a permutation/exact test where it does not.
Correctness is \emph{always} reported beside time: a condition that is faster but
less correct is not ``better.'' Speed without correctness is meaningless.

\paragraph{Scoring procedure and reliability (pre-specified).} Correctness is the
one primary measure produced by human judgement rather than instrumentation, so
the procedure is fixed in advance. (a) \emph{Blinding.} Scoring works from a
final-state export --- board state, posted messages, PR-review verdict and
comment --- with \code{agent\_condition}, the agent transcript, and all timing
stripped. Knowing a participant was in C3 while judging their output is a direct
route to confirming H1, and the export makes blinding cheap, so it is mandatory
rather than aspirational. (b) \emph{Two independent coders} score every task that enters an analysis, with
no sampling shortcut and nothing single-coded.

\emph{Scope of the correctness coding: the agent round only.}\label{sec:correctness-scope} Every correctness
model in this plan fits on agent-phase tasks only; the six baseline scores per
participant would feed the balance check (\S\ref{sec:balance}) and the $\kappa$
sample, and nothing else. We therefore fix, in advance, that the coded unit set
is the $6\times N$ \emph{agent} tasks, and the baseline round is not run unless a
reviewer asks for it. This is a scoping decision about the \emph{unit set}, not a
sampling shortcut, and the distinction is what makes it safe: whichever tasks are
coded are coded in full by both coders, and $\kappa$ is reported over the whole
coded set rather than a subset of it. What is given up is small and is replaced
rather than lost --- the balance check still runs on baseline time and baseline
effort, both instrumented for every participant, and baseline \emph{time} is the
covariate the primary model actually carries (Eq.~\eqref{eq:primary}), so the arm
comparability question is answered by the measure that matters with or without
the coded half. What is bought is a halving of the largest of three codings. The
raw final-state export is retained for all twelve tasks and
\code{coding\_packet.py --phase baseline} builds the second round on demand, so
the decision is reversible in the only direction that could matter. (c) \emph{Reliability} is reported as quadratic-weighted
Cohen's $\kappa$ on the 0--3 scale (weighted because the score is ordinal --- a
3-vs-2 disagreement is not a 3-vs-0 disagreement), with exact agreement rate
alongside it. (d) \emph{Disagreements} are resolved by discussion to consensus,
still blind to condition; the pre-adjudication $\kappa$ is what gets reported,
not the post-consensus one. (e) The rubric is fixed and piloted on the pilot
participants' data before any real session is scored, and is not revised once
scoring of the real sample has begun.

\subsection{The coded set is capped by enrolment order}\label{sec:cap}

Recruitment and coding scale differently, and the design has to say so before
recruitment opens rather than after. Sessions are unproctored and self-serve
(\S\ref{sec:timing}), so one more participant costs the investigators nothing and
improves every instrumented measure --- including the \emph{primary} outcome,
baseline-adjusted task time --- for free. The three human codings cost roughly an
hour and a half of two people's time per participant, are linear in $N$, and
constitute essentially the entire investigator cost of the study. Recruitment is
by open link through course lists and developer communities, so $N$ may
comfortably exceed the 20--30 target. Capping \emph{recruitment} to protect the
coding budget would discard the free half of the study to preserve the expensive
half, and would shrink the sample on exactly the contrast this design is least
able to estimate.

So we cap the \emph{coded} set instead, by a rule fixed here in advance: the
first \textbf{ten participants per condition, in enrolment order}
(\code{provisioned\_at}, ties broken by participant code) enter the human-coded
set; everyone else contributes to every instrumented analysis and to none of the
hand-coded ones. Eligibility is having produced at least one agent-phase final
state --- presence of data, not quality of it. The same set is used for all three
codings, so correctness, agent error and themes describe one group of people
rather than three. The rule is implemented in
\code{coding\_packet.coded\_set}, which writes the full manifest --- selected and
unselected alike, with the reason --- so the decision is auditable rather than
implicit in which files exist.

\paragraph{Why this is a cap and not a choice of sample.} Three properties, and
each is doing distinct work.

\emph{Enrolment order is fixed before the session happens.} It is stamped at
provisioning, when the participant has done nothing yet, so no property of how a
session went can move anyone into or out of the coded set. Every alternative we
considered fails this test in the same way: selecting completers excludes on
dropout, which correlates with task difficulty; selecting on data quality or on
having used the agent selects on the dependent variable outright. Those would be
indistinguishable from choosing the sample after seeing it. This one is not,
because there is nothing to see at the moment the rule applies.

\emph{The cap is applied within condition, and condition is not correlated with
enrolment order.} Condition assignment is adaptive minimisation --- each
participant goes to the currently least-filled arm
(\code{provision.rs::pick\_agent\_condition}) --- which is already chosen to break
the confound between condition and enrolment order (\S\ref{sec:limits}). Taking a
prefix of that order therefore cannot load one arm with early enrollees and
another with late ones, and applying the cap per arm guarantees equal coded cells
even if recruitment itself runs uneven.

\emph{The residual risk is drift, and we test it rather than assert it.} People
who sign up in the first week may differ from those who sign up in the third:
keener, or reached through a different channel. A prefix of enrolment order
inherits that difference. Two things follow, both pre-specified. First, it does
not bias the between-condition contrasts --- the cap is applied within condition,
so any drift lands on all three arms alike, which is the same argument
randomisation makes about any nuisance variable. Second, it does bound how far
the correctness result generalises to the full collected sample, and that bound
is measured: \code{analyze.coded\_set\_comparison} compares coded against uncoded
participants on baseline time, baseline effort, agent-task time and agent
invocation rate (Mann--Whitney), all of which are instrumented and therefore
exist for uncoded participants precisely because the cap removes hand coding and
not telemetry. No difference is evidence the coded set stands in for the
collected one; a difference is reported with the coverage table and carried as a
generalisation limit.

\paragraph{The asymmetry is reported, not smoothed over.} The consequence of an
uncapped $N$ and a capped coded set is that time is estimated on more people than
correctness is. This plan's standing commitment is that correctness is
\emph{always} reported beside time (\S\ref{sec:primary}), and a speed result from
$N$ participants set beside an accuracy result from fewer is not that pairing. So
two things are produced on every run. \code{tables/coding\_coverage.csv} gives
collected, coded, and items scored per condition, and it is reported with the
participant flow rather than in a limitations paragraph. And the sensitivity
panel gains a fifth row --- the confirmatory contrast refitted on the
correctness-coded participants alone --- so the headline and the correctness
result can be read on a common set of people. That row is a \emph{variant}, not
the model of record: discarding collected, instrumented observations to match the
coding budget would be the wrong direction, since the extra participants are
precisely what buys precision on the primary outcome. But a headline that exists
only in the uncoded remainder would be visible immediately, which is the point.

\subsection{Cognitive load (H3) --- two measures, two analyses}

NASA-TLX is collected \emph{twice}---once in a post-baseline interstitial (manual
phase, \code{tlx\_*\_baseline}) and once at exit (agent phase, \code{tlx\_*})---so
it now yields a within-person, phase-to-phase delta. We test H3 two ways:
\begin{enumerate}[nosep,label=(\alph*)]
  \item \emph{global TLX delta} (primary): per participant,
  $\Delta\textit{TLX} = \textit{TLX}_{\text{agent}} - \textit{TLX}_{\text{baseline}}$
  on the composite (and per subscale)---%
  \footnote{\code{tlx\_performance} (and \code{tlx\_performance\_baseline}) is
  anchored Failure$\rightarrow$Perfect, i.e.\ high $=$ good performance, the
  opposite of the other five items. Reverse-score it once
  ($\text{MAX}+\text{MIN}-\text{raw}$) so higher $=$ more load before it enters
  the composite; leave the other five as stored.}
  then compare $\Delta\textit{TLX}$
  \emph{between} conditions (Welch ANOVA / Kruskal--Wallis at small $N$). H3
  predicts a larger drop (more negative $\Delta$) under broader context. The
  paired baseline removes the between-person level differences that a single
  end-of-study TLX confounds;
  \item \emph{per-task effort}: $\textit{effort} \sim \textit{condition} \times
  \textit{task\_type} + \textit{baseline\_effort} + (1\mid\textit{participant})$,
  carrying each participant's type-matched baseline-task effort as a covariate
  (the same ANCOVA logic as Eq.~\eqref{eq:primary}; the finer-grained
  within-person test).
\end{enumerate}
Both are reported; (a) is the headline load result, (b) localizes it by task type.

\subsection{Perception--behavior gap (H5) --- the most novel result}

\emph{Across} participants (each participant is one point---one exit-survey
trust value against one summary performance value), Spearman-correlate
\code{trust\_speed\_perceived} and a trust composite against the \emph{measured}
performance. The single trust item is correlated against a per-participant
summary of normalized time and correctness: pre-specified as the mean across the
two cross-surface agent tasks (T5, T6)---the tasks where the agent's measured
effect is largest, so the perceived--actual comparison is sharpest---with the
all-six-agent-task mean reported as a robustness check. We then characterize
divergence: people tend to trust an agent that visibly \emph{does
more} (C3 acts across all surfaces) even when it was not actually faster, or was
occasionally wrong. We also cross \code{trust\_rely} (willing to rely) against
\code{trust\_verify} (felt the need to double-check) for over-/under-reliance,
and pair \code{int\_switch\_perceived} (perceived switching) with the behavioral
view-switch rate --- the one place a self-report is checked against a behavioral
counterpart of the same construct rather than against an outcome.

\paragraph{Multiplicity (pre-specified).} Crossing five perception measures with
four performance summaries, pooled and within each condition, is roughly a
hundred correlations; at $\alpha = 0.05$ five would clear significance with no
effect present at all, and the per-condition rows are $n = 7$--$10$, where one
participant moves $\rho$ by $0.3$. So exactly one correlation is fixed in advance
as primary --- \code{trust\_speed\_perceived} against mean normalized time on the
cross-surface agent tasks, pooled --- and reported unadjusted. Every other cell is
exploratory, carries a Benjamini--Hochberg adjusted $p$ over that family, and the
count of tests is reported with the table. We read none of them as findings on
their own. Separately, if Cronbach's $\alpha$ for the three-item trust composite
falls below the pre-committed $0.70$, every composite row is demoted to
descriptive and the three items are reported individually.

\paragraph{The integration-perception items.} \code{int\_findacross},
\code{int\_noticed}, \code{int\_vs\_separate}, \code{int\_reuse} and
\code{int\_switch\_perceived} are reported per condition as descriptives. No
hypothesis rides on them --- scope is manipulated, perceived integration is not,
and they are not a validated scale, so they are not combined into a composite.
They are reported because a reader is entitled to ask whether participants
noticed the thing the paper is about, and because a flat profile across
conditions is itself informative about how visible agent scope was.

\subsection{The single cleanest test: task T6}\label{sec:t6}

T6 (link the PR to a card and note it in chat) is constructed so the correct
pull request is identifiable \emph{only} with GitHub context: chat refers to it
vaguely; the PR's title/body disambiguates. Therefore: if C3 beats C2 on T6, the
top rung earns its place; if C3 does \emph{not} beat C2 on T6, the data
themselves say GitHub context is not pulling its weight. One task carries a large
share of the argument, in either direction.

\section{Reliance and division of labor (RQ2)}

RQ2 asks how users actually work with the agent. Alongside the primary outcome
models, the behavioral telemetry gives a division-of-labor picture --- does the
user delegate to the agent or take over, and do they check its output?

\paragraph{H4's calibration half needs its own ground truth.} ``Delegation rises
with scope'' is the delegation ratio below. ``Reliance is imperfectly calibrated
to where the agent is actually correct'' is a different claim and cannot be
tested from the same data: it requires a judgement about whether \emph{the
agent} was right, which is not the same as whether the participant's task ended
right. A participant who spots a wrong suggestion and fixes it produces a fully
correct outcome from a wrong agent --- exactly H4's case, and exactly the case an
outcome-based measure cannot see. We therefore code agent-error type per task
(0 correct / 1 partially wrong / 2 wrong) from the exported transcript
(\code{study\_agent\_turns}), and cross it against override behavior
(\code{agent\_suggestion\_rejected}, \code{agent\_output\_edited},
\code{agent\_action\_undone}) on tasks where the agent was invoked. The quantity
is the override rate within each stratum; good calibration means overriding far
more often when the agent was wrong, and H4 predicts a gap smaller than that. A
gap near zero is the strong form of the claim, not a null result.

\paragraph{What is blinded here, and what is not.} The transcript gives the
condition away --- a C1 agent never acts, a C3 one discusses pull requests --- so
this coding cannot be condition-blind, and stripping the transcript would remove
the thing being coded. That matters less than it first appears: condition is
\emph{not} the contrast in the calibration test, which compares override rates
between strata of agent correctness, so a coder who infers the condition has no
route to move the result. It is not nothing, and we report the coding as
unblinded to condition.

The variable that must be hidden is the \emph{outcome}. A coder who reads a
participant writing ``no, that's the wrong PR, undo it'' and concludes the agent
was wrong has inferred the stratifying variable from the very behavior the
hypothesis predicts; the test would then confirm itself out of nothing. So the
packet truncates each exchange after the agent's last turn, and the residual ---
an intermediate participant turn can also react --- is stated rather than solved.
Coders are directed to settle factual claims against a supplied answer key
rather than against the participant.

\paragraph{Two error codes, because one of them is checkable.} Every participant
works in a clone of one fixed seed (12 cards, 18 messages, 3 issues, 3 pull
requests), so most claims an agent makes have a lookup answer: whether PR 34 is
the vaccination-reminder one is a fact, not an impression. We therefore code two
independent binaries: \emph{factual error} (the workspace contradicts what the
agent said, settled against the answer key) and \emph{judgement error}
(factually fine but the wrong move for the task). \textbf{Factual error is the
pre-specified stratifier}, because it is the closest available reading of ``where
the agent is actually correct'' and because two coders can agree on it;
judgement error is reported separately as the softer measure. Agreement is
reported for each separately --- pooling them would hide that one is a lookup and
the other an opinion.

\paragraph{A hypothesis that a null would confirm (pre-specified reading).} H4's
calibration clause predicts \emph{imperfect} calibration, i.e.\ a
\emph{small} gap between the override rate when the agent was wrong and when it
was right. That is a hypothesis satisfied by a null result, which at
$n \approx 8$ per cell is exactly the shape that invites a wide, uninformative
interval to be written up as confirmation. We fix the reading in advance: the
gap is reported as a difference in proportions \emph{with its confidence
interval}, and an interval spanning both zero and a substantial gap is reported
as \textbf{uninformative about H4}, not as support for it. This is the same
commitment already made for C2-vs-C3 (\S\ref{sec:primary}) and it is made here
for the same reason. Only a gap estimated with enough precision to exclude
good calibration counts as evidence for H4.

To make that operational rather than rhetorical, we fix the comparison value
now: \textbf{a 25-percentage-point} higher override rate on wrong-agent tasks
than on right-agent tasks is taken as the standard of good calibration
(\code{analyze.CALIBRATED\_GAP\_PCT}). It is deliberately unambitious --- H4
predicts users fall short of it --- and it is fixed in advance precisely because
a threshold chosen after seeing the estimate could be moved to make any result
support the hypothesis. The interval is computed by Newcombe's method for a
difference of two proportions, which stays inside $[-1,1]$ at the small
per-stratum counts this design produces where a normal approximation does not.
The threshold is used only to classify the interval as informative or not; it
is never a test, and no result is declared by comparing a point estimate to it.

Correctness stays in its own separate blinded packet and is coded first, so the
load-bearing measure is never contaminated by this one. Where the agent-error
coding is unavailable the analysis falls back to outcome correctness and says
so, labeled the weaker measure it is: a task can end correct \emph{because} the
participant fixed a wrong suggestion, so the fallback under-detects exactly the
case H4 is about.

\paragraph{Adjudication.} Disagreements on all three codings (correctness,
agent error, themes) are resolved by an independent third coder who sees the
disputed items and neither coder's choice, rather than by discussion to
consensus. Discussion has a known failure mode at this scale --- the more
confident or more senior coder tends to win, and where that person knows the
hypotheses, bias re-enters at precisely the step blinding existed to prevent.
The third coder is shown only the items in dispute, so the cost is proportional
to disagreement rather than to the corpus. Reported $\kappa$ remains the
pre-adjudication figure in every case.

Supporting telemetry descriptives:
\begin{itemize}[nosep]
  \item \emph{delegation ratio} by condition (agent-executed vs.\ user-performed
  actions)---%
  \footnote{The chat tool \emph{drafts}: the agent never posts, so the
  participant's send is recorded as an ordinary user action. A drafted message is
  therefore excluded from the agent-executed numerator---it is flagged
  \code{drafted} on \code{agent\_tool\_call}---because counting it would put one
  outcome in the numerator and the denominator at once. Draft counts are reported
  separately so that channel of delegation stays visible. C1's ratio is $0$ by
  construction rather than by participant choice, so C1's reliance signal is
  invocation and advice-following, not this ratio.};
  \item the \emph{agent funnel}: \code{agent\_invoked} $\rightarrow$ tool attempted $\rightarrow$ \code{executed} $\rightarrow$ \code{accepted}/\code{edited}/\code{undone};
  \item tool error / rejection rates; cross-reference traversal counts; view-switch rate;
  \item a CONSORT-style participant flow from \code{study\_participants.status} $+$ \code{session\_abandoned};
  \item thematic analysis of the four open-ended items (\code{open\_useful},
  \code{open\_frustrating}, \code{open\_change}, \code{open\_integration}).
  Procedure, fixed in advance so it is not improvised on results day: one coder
  reads every response and revises a starter codebook, the codebook is then
  frozen, and the coders assign \emph{one primary theme} per
  response. The single-theme constraint is what makes Cohen's $\kappa$
  well-defined; secondary impressions go in an unscored notes column. $\kappa$ is
  unweighted --- themes are nominal, and a weighted $\kappa$ would invent an
  ordering the codebook does not have --- and, as with correctness, is reported
  pre-adjudication.

  \emph{This is the one coding that is not fully double-coded, and it follows
  from the status this plan already assigns it.} Coder 1 codes every response;
  coder 2 independently codes a stratified \textbf{25\,\%} reliability sample,
  drawn across all four items so no question goes untested, and $\kappa$ is
  computed on that overlap alone. Three reasons this is the right place to spend
  the saving, and nowhere else. It is the only one of the three codings whose
  output is not a number in a model: the themes are labelled illustrative rather
  than confirmatory two sentences below, and a double-coded reliability sample
  with a reported $\kappa$ is the standard instrument for exactly that status,
  where full double-coding is the standard for a measure carrying a hypothesis.
  It is the coding that grows fastest with $N$ --- four responses per participant
  against six tasks --- so it is where linear growth bites first. And the
  reliability sample is not small in absolute terms: at $N$ past the cap it is
  around forty-five double-coded responses, more than many published thematic
  analyses report $\kappa$ over. Correctness and agent error stay fully
  double-coded, with no sampling shortcut, because they feed pre-specified tests.
  The saving is made visible rather than taken quietly: the packet tells each
  coder which role they hold and how large the overlap is, and the reported
  $\kappa$ carries the overlap size beside the corpus size, so a reader can never
  mistake it for agreement over the whole set. Single-coded responses take coder
  1's theme as their value --- there is no disagreement to adjudicate --- and are
  excluded from $\kappa$ by construction; filling them in and \emph{then}
  computing agreement would report a figure inflated by rows one person saw
  alone. Responses are shuffled across participants and items and
  carry no condition label; the residual risk that an answer reveals its own
  condition (``it linked the PR for me'') is stated rather than papered over,
  and the qualitative findings are reported as illustrative, not confirmatory.
\end{itemize}

\section{Figures}

Five figures do the heavy lifting. (Placeholders below; swap for real plots after
data collection.)

\begin{figure}[h]\centering
\figplaceholder{\textbf{Fig.\ 1 --- Normalized time by condition $\times$ task
type.} Grouped bars/lines: task type on the x-axis (comprehension /
single-surface / cross-surface), series C1/C2/C3. Expected shape: roughly equal
on the left (control tasks), fanning apart on the right with C3 lowest. This
picture \emph{is} H2.}
\caption{The headline interaction plot.}
\label{fig:interaction}
\end{figure}

\begin{figure}[h]\centering
\figplaceholder{\textbf{Fig.\ 2 --- Correctness by condition $\times$ task type.}
Same axes as Fig.\ 1, shown adjacent so speed and accuracy are read together---
evidence that any speed gains are not from cutting corners.}
\caption{Correctness, paired with the timing figure.}
\label{fig:correctness}
\end{figure}

\begin{figure}[h]\centering
\figplaceholder{\textbf{Fig.\ 3 --- Perception vs.\ behavior.} Scatter; each dot a
participant. X: measured speedup (normalized time). Y: perceived speedup
(\code{trust\_speed\_perceived}). Departure from the diagonal \emph{is} the H5
story.}
\caption{The perception--behavior gap.}
\label{fig:percep}
\end{figure}

\begin{figure}[h]\centering
\figplaceholder{\textbf{Fig.\ 4 --- Delegation ratio by condition.} How much the
agent did vs.\ the user did, by condition---work shifting off the user as context
broadens (RQ2).}
\caption{Division of labor across the ladder.}
\label{fig:delegation}
\end{figure}

Figure~5 is the scoping manipulation-check table (\S\ref{sec:manip}), not a plot.

\section{What comes first in Results}

Two things precede every inferential result.

\subsection{Manipulation check (scoping held)}\label{sec:manip}
Before claiming ``C1 had less context,'' we prove it from
\code{agent\_context\_sent}: 100\,\% of C1 turns contained no board/GitHub data,
and 100\,\% of C2 contained no GitHub. This converts the independent variable
from an assumption into evidence. It is a small table, but load-bearing.

\subsection{Baseline balance (randomization check)}\label{sec:balance}

Condition is assigned at provisioning, but the six baseline tasks are identical
across arms and run with no agent, so baseline performance is a check on whether
the arms were comparable \emph{before} the manipulation existed. Minimization
keeps the cells equal in size; it cannot make them equal in ability, and at
$n \approx 8$ per cell a chance imbalance is entirely possible. We compare
baseline time and baseline effort across conditions (one value per participant,
Kruskal--Wallis). The expected result is no
difference. A difference is not a finding about the manipulation --- it is a
caveat on every contrast that follows, and it is reported with the flow numbers
rather than in a limitations paragraph at the end.

Baseline \emph{correctness} is absent from that list because the baseline
correctness round is not run (\S\ref{sec:correctness-scope}). The check survives
the omission on its own terms: both remaining measures are instrumented for every
participant --- so this check, unlike the correctness models, runs on the full
collected sample rather than the coded subset --- and baseline \emph{time} is the
one the primary model actually carries as a covariate
(Eq.~\eqref{eq:primary}), so the imbalance that would propagate into the
contrasts is exactly the imbalance still being tested for. If the baseline round
is coded later, correctness rejoins the table and Results says so.

\subsection{Participant flow and timing validity}\label{sec:timing}
A CONSORT-style count (started / finished / dropped, from
\code{study\_participants.status} and \code{session\_abandoned}), plus a statement
that task time is \emph{active engaged time}: the client clock runs only while
the tab is visible, pausing on background/close and resuming on return, and
persists across reloads/re-entry, so away-time is excluded by construction rather
than after the fact. Each task also logs \code{away\_ms}, \code{pause\_count},
and \code{was\_interrupted}. Pre-specified handling: analyze active time as
primary; run a sensitivity analysis excluding \code{was\_interrupted} tasks;
report the interruption rate.

\paragraph{One sensitivity panel, computed every run.} The five fits the plan
commits to --- primary (intention-to-treat, active time), idle-adjusted,
interrupted tasks excluded, as-used, and the coded subset
(\S\ref{sec:cap}) --- are all produced on every run and the
confirmatory contrast is reported from each, side by side, with the number of
observations and participants each was fitted on. Which one is of record is fixed
in advance (the primary row); the rest exist so that a headline which moves
across them is reported as moving. Choosing among them after seeing the numbers
is how a pre-specified sensitivity quietly becomes a post-hoc choice, and
producing all four unconditionally is what makes that impossible. The idle rule's
trigger count and the interruption rate are reported per condition, not just
pooled: a correction that lands unevenly across arms is a confound rather than a
cleanup, and that is the thing a reader needs to be able to check.

\paragraph{Visible-idle rule (pre-specified).} The clock pauses when the tab is
hidden, which catches the common way of stepping away but not all of it: a
participant who leaves the tab in front and walks off is still banking active
time. Because sessions are unproctored and $n\approx8$ per cell, one such task
can carry a condition. We therefore pre-specify, before any data are seen: within
a task, a stretch of $\geq$\,180\,s of visible time containing \emph{zero}
telemetry events is treated as idle and subtracted from active time. The
threshold sits well above the natural gap for reading a channel or a diff
(tasks run $\sim$4--5\,min in total and emit view, card, message, and agent
events throughout) and well below the scale of a genuine walk-away. We report how
many tasks trigger the rule, and give the primary model both with and without the
correction; a result that depends on which is used will be reported as such
rather than resolved in the direction we prefer.

\paragraph{Practice effects bound what the baseline can claim.} Phase order is
fixed --- the six manual baselines always precede the six agent tasks, because
seeing the agent cannot be undone. So by T1 every participant is six tasks more
fluent in Weld, and the agent$-$baseline delta confounds agent help with ordinary
learning. This does not threaten the study's inferences, which are all
\emph{between} conditions where the ordering is identical by construction. It
does mean one sentence is unavailable to us: we never claim that the agent made
participants faster than working manually. H1 is stated and tested as a
difference in the baseline-adjusted agent-task time \emph{across conditions}
(Eq.~\eqref{eq:primary}), never as an absolute improvement over C0.

\section{Reporting order (Results section skeleton)}

\begin{enumerate}[nosep]
  \item Sample \& flow --- who, how many, dropouts (CONSORT-style, per condition).
  \item Coding coverage --- collected vs.\ hand-coded per condition, and the coded-vs-uncoded comparison (\S\ref{sec:cap}). Reported here, with the flow, rather than in Limitations: a reader meets the two sample sizes before meeting any result computed on them.
  \item Baseline balance --- were the arms comparable before the agent existed (\S\ref{sec:balance}).
  \item Manipulation check --- scoping held (\S\ref{sec:manip}).
  \item RQ1 / H1--H2 --- normalized time and correctness by condition $\times$ task type (Figs.\ \ref{fig:interaction}, \ref{fig:correctness}); the core result. The confirmatory contrast from the model, then the sensitivity panel and agent uptake.
  \item H3 --- cognitive load (between-condition TLX delta, agent vs.\ baseline phase; per-task effort ANCOVA within person).
  \item H5 --- perception vs.\ behavior (Fig.\ \ref{fig:percep}); the memorable result. Primary correlation first, exploratory family after, with its test count.
  \item RQ2 --- reliance and division of labor (Fig.\ \ref{fig:delegation}), the H4 calibration cross-tab, open-ended themes.
\end{enumerate}

Every number above is produced by \code{paper/analysis/}, one table per file,
by \code{make run}. The report is written in this order deliberately, so that
assembling the Results section is transcription rather than re-analysis.

\section{Hypothesis-to-analysis map}

\begin{table}[h]
\centering\small
\begin{tabularx}{\linewidth}{@{}l L L@{}}
\toprule
\textbf{H} & \textbf{Test} & \textbf{Key fields} \\
\midrule
H1 & condition effect on log active time, baseline-adjusted ANCOVA, Eq.~\eqref{eq:primary}; confirmatory C3-vs-C1 cross-surface contrast taken from the fitted model, with the ratio contrast reported beside it as the descriptive companion & \code{task\_time\_ms}, C0 baseline (covariate), \code{agent\_condition} \\
H2 & \code{condition $\times$ task\_type} interaction, Eq.~\eqref{eq:primary} (supporting) & \code{task\_type} across the 6 agent tasks \\
H3 & (a) between-condition TLX delta (agent$-$baseline); (b) per-task effort ANCOVA on the type-matched baseline effort & \code{tlx\_*\_baseline}, \code{tlx\_*}, \code{task\_rating.effort} (baselines included) \\
H4 & (a) delegation ratio by condition; (b) override rate crossed with coded agent-error type, on invoked tasks & \code{agent\_tool\_call} (executed), \code{agent\_suggestion\_*}, \code{agent\_output\_edited}, \code{agent\_action\_undone}; \code{study\_agent\_turns} $\rightarrow$ agent-error coding \\
H5 & trust vs.\ measured speed/correctness (Spearman); one pre-specified primary correlation, remainder FDR-adjusted & \code{trust\_speed\_perceived}, trust block, normalized time \\
--- & scoping manipulation check & \code{agent\_context\_sent} \\
--- & participant flow; timing-validity sensitivity panel (ITT / idle-adjusted / interruptions excluded / as-used) & \code{study\_participants.status}, \code{session\_abandoned}, \code{task\_skipped}, \code{was\_interrupted}, \code{agent\_invoked} \\
\bottomrule
\end{tabularx}
\caption{Each hypothesis, its test, and the fields it draws on.}
\label{tab:hmap}
\end{table}

\section{Limitations and small-$N$ strategy}\label{sec:limits}

Three between-subjects cells at $N\approx20$--$30$ ($\sim$7--10 per cell) is a
modest, focused sample---a short full study, not a large trial, so we report
estimates with their uncertainty rather than claiming a definitively powered
omnibus interaction. What makes it viable: the within-person baseline (absorbed
as a covariate, Eq.~\eqref{eq:primary}) and six agent observations per
participant (more observations than people).
Mitigations: condition assignment is balanced by adaptive minimization---each
participant is assigned to the currently least-filled condition, ties broken at
random (implemented in \code{provision.rs::pick\_agent\_condition}), which keeps
the cells near-equal and breaks the confound between condition and enrollment
order; this plan itself, which fixes the confirmatory focused contrast
(\S\ref{sec:t6}) and the exploratory remainder in advance and is archived with the
replication package rather than deposited in a public registry; and reporting
effect sizes and confidence
intervals---and, where appropriate, permutation or Bayesian estimation---rather
than relying on null-hypothesis significance tests. We state plainly that this is
a focused study, not a large trial.

\paragraph{Two sample sizes, and which claim rests on which.} Because
recruitment is uncapped and the coded set is not (\S\ref{sec:cap}), this study
reports on two nested samples, and we would rather name the consequence than
have a reader find it. Everything instrumented --- the primary timing model and
its confirmatory contrast, TLX, effort, the survey, delegation, the funnel, the
manipulation check --- rests on the full collected sample. Everything requiring
human judgement --- correctness, agent error, themes --- rests on the coded
subset, which is the first ten participants per condition in enrolment order.
The limitation this creates is one of \emph{precision on the judged measures},
not of validity: the coded subset is selected by a rule fixed before any session
ran and applied within condition, so it cannot have been chosen to favour a
result, and \code{coded\_set\_comparison} reports whether coded and uncoded
participants differ on the instrumented measures at all. The honest reading, and
the one we commit to: if the correctness result is imprecise, it is imprecise
because we could afford to code thirty people and not because we recruited
thirty. Coding more is a resourcing decision that can be revisited without
re-running the study, since the raw final-state export is retained for every
participant. Nothing about the cap is recoverable only by collecting again ---
which is exactly why the cap goes here and not on recruitment.

\paragraph{The top rung is a narrow rung.} Under study mode GitHub is served from
fixtures and is read-only: pull requests carry one or two files of diff, and the
write endpoints are locked. C3's capability beyond C2 therefore amounts to
reading small diffs and linking a PR to a card. This bounds the size of any
C2$\rightarrow$C3 effect we could detect, and it bounds generalization: in a real
repository, review is the expensive part of the work, and that is precisely where
GitHub-scoped agents might earn or lose their keep. The lockdown is deliberate
(no participant action can reach a real repository) and we would make the same
choice again, but the resulting claim is about a GitHub rung of modest depth, not
about code review at realistic scale.

\paragraph{One model, one moment.} The agent is a single pinned model snapshot,
fixed temperature, held constant for the whole run and named exactly in the
methods. That buys internal validity --- no mid-study capability drift --- at the
cost of external validity: results describe how users worked with \emph{this}
agent, and a materially stronger or weaker model could move both effectiveness
and the trust calibration H5 measures. The snapshot is also chosen to be capable
enough that the top rung is not floored by the model rather than by the design;
a null under a weak agent would be uninterpretable.

\paragraph{Availability, not instruction.} Agent-task wording tells participants
the assistant \emph{is available}, never to use it, precisely so RQ2 measures
chosen delegation rather than compliance. The cost is exposure: a participant may
complete agent tasks with little or no agent use, which is informative for RQ2
but reduces the effective sample for the effectiveness contrasts. We report agent
uptake per condition alongside the flow numbers, and pre-specify that
effectiveness models are fit on all tasks as assigned (intention-to-treat), with a
secondary as-used analysis restricted to tasks with at least one
\code{agent\_invoked}; divergence between the two is reported, not resolved.

\end{document}
